\documentclass[12pt,hyperref,a4paper,UTF8]{ctexart}
\usepackage{HDUReport}
\usepackage{listings}
\usepackage{xcolor}
\usepackage{graphicx}
\usepackage{setspace}
\usepackage{float}
\setstretch{1.5} % 设置全局行距为1.5倍

\usepackage{enumitem} % 载入enumitem包以便自定义列表环境
\setlist[itemize]{itemsep=0pt, parsep=0pt} % 设置itemize环境的项目间距和段落间距

\setmainfont{Times New Roman} % 英文正文为Times New Roman


\usepackage{tikz}
\usetikzlibrary{shapes.geometric, arrows}
\usetikzlibrary{positioning, arrows.meta}
\usetikzlibrary{calc}
%封面页设置
{   
    %标题
    \title{ 
        \vspace{1cm}
        \heiti \Huge \textbf{《电磁波与电磁场实验》作业报告} \par
        \vspace{1cm} 
        \heiti \Large {\underline{思政报告}   }     
    }

    \author{
        \kaishu\Large 指导教师\ \dlmu[9cm]{吴爱婷}  \\
        \kaishu\Large 学生姓名\ \dlmu[9cm]{陈文轩} \\ %学院
        \kaishu\Large 学生学号\ \dlmu[9cm]{23040447} \\ %
        \kaishu\Large 学生班级\ \dlmu[9cm]{23186011} \\ %班级
        \kaishu\Large 学生专业\ \dlmu[9cm]{智能硬件与系统(电子信息工程)} \qquad  \\ %学号
        \kaishu\Large 实验日期\ \dlmu[9cm]{2025.10.20} \qquad \\
         %姓名 
    }
        
    \date{\today} % 默认为今天的日期，可以注释掉不显示日期
}
%%------------------------document环境开始------------------------%%
\begin{document}

%%-----------------------封面--------------------%%
\cover
\thispagestyle{empty} % 首页不显示页码
%%------------------摘要-------------%%
%\newpage
%\begin{abstract}




%\end{abstract}

%\thispagestyle{empty} % 首页不显示页码

%%--------------------------目录页------------------------%%
% \newpage
% \tableofcontents
% \thispagestyle{empty} % 目录不显示页码

%%------------------------正文页从这里开始-------------------%
\newpage
\setcounter{page}{1} % 让页码从正文开始编号

%%可选择这里也放一个标题
%\begin{center}
%    \title{ \Huge \textbf{{标题}}}
%\end{center}

\section{电磁场与微波技术对于社会发展的重要性}

电磁场与微波技术是现代科技的重要组成部分，对社会发展具有深远的影响。首先，在通信领域，微波技术支撑了无线通信、卫星通信和移动网络的发展，使得全球信息交流更加便捷高效。例如，5G和即将到来的6G技术均依赖于微波频段的高频特性。其次，在医疗领域，微波技术应用于肿瘤治疗、成像诊断和物理疗法，如微波消融手术和MRI设备。此外，微波技术在工业、航空航天和国防等领域发挥关键作用，如雷达系统、卫星导航和电子战装备。最后，微波技术促进了物联网和智能家居的发展，推动了数字化转型和社会进步。然而，随着技术的广泛应用，也带来了电磁辐射对人体健康的潜在影响和社会安全问题，需要加强监管和安全设计。

\section{电磁场与微波技术相关法律法规及安全健康的微波电路设计}

电磁场与微波技术相关法律法规主要包括国际和国内标准，以确保电磁辐射的安全使用。在国际上，国际非电离辐射防护委员会（ICNIRP）制定了电磁场暴露限值指南，限制电磁辐射对人体的暴露量。在中国，国家标准GB 8702-2014《电磁环境控制限值》规定了电磁辐射的控制标准，适用于微波设备的设计和使用。此外，《中华人民共和国电磁辐射环境保护法》要求设备制造商和使用者遵守辐射限值，避免对公众健康的危害。

在微波电路设计中，实现安全健康的设计需要遵循以下原则：首先，进行电磁兼容性（EMC）分析，确保电路不会产生过多的电磁干扰。其次，采用屏蔽技术和滤波器减少辐射泄漏。第三，使用低功率设计和自动功率控制（APC）机制，避免不必要的辐射。最后，进行人体暴露模拟和测试，验证设计符合安全标准。通过这些措施，可以设计出既高效又安全的微波电路，保护用户健康和社会环境。

\section{国内外微波设计仿真软件的现状和特点}

国内外微波设计仿真软件的发展现状各有特点。国外软件处于领先地位，主要代表包括Keysight公司的Advanced Design System (ADS)、Ansys公司的HFSS和Dassault公司的CST Studio Suite。这些软件具有强大的电磁场模拟能力，支持多物理场耦合分析、优化设计和虚拟原型验证，广泛应用于学术研究和工业设计。其特点是算法成熟、精度高、用户界面友好，但价格昂贵，且核心技术受限于国外供应商。

国内微波设计仿真软件起步较晚，但近年来取得显著进步。例如，中国电子科技集团开发的软件如Microwave Office的本土化版本，以及北京大学等高校开发的开源工具如OpenEMS。这些软件在特定领域如毫米波和太赫兹仿真方面表现出色，但整体上与国外软件相比，在算法深度、计算效率和生态系统方面存在差距。国内软件的特点是成本较低、易于定制，但功能覆盖面窄、稳定性有待提升。

正视国内外技术差距，我们应树立自力更生的思想。通过加大研发投入、人才培养和国际合作，缩小差距，实现微波设计软件的自主可控。这不仅有助于提升国家科技安全，也能推动相关产业的健康发展。

\section{电磁场与微波技术相关产品的自主创新设计}

实现电磁场与微波技术相关产品的自主创新设计，需要从多个方面入手。首先，在设计仿真软件方面，积极采用国产软件如Microwave Office本土化版本或自主开发的工具，进行电路和系统的模拟验证。这有助于减少对国外软件的依赖，提升设计效率。其次，推动元器件国产化，使用国产芯片、滤波器和放大器等组件，降低供应链风险。例如，支持国内厂商如华为海思或中兴微电子的元器件应用。第三，发展国产材料，如高性能陶瓷基板和复合材料，用于微波电路的制造，确保材料性能和供应稳定。最后，坚持独立自主设计理念，从需求分析到原型测试，全程由国内团队主导，避免技术泄露。通过这些措施，可以实现产品的自主创新，推动电磁场与微波技术的可持续发展。

\end{document}